% !TeX program = lualatex
% ================================================================
% Urdu Tier 3 Vocabulary Version of ScalingVolumeSurfaceAreaStarter
%
% Generated by the server using the following settings:
%
%   language            Urdu (urd)
%   text direction      right to left
%   font                Noto Nastaliq Urdu
%   fallback font       Noto Serif
%   built from          latex/starters/targeted/KS4/geometry/scalingVolumeSurfaceAreaStarter.tex
%   vocabulary key      latex/starters/targeted/KS4/geometry/scalingVolumeSurfaceAreaStarter/L2/urd/scalingVolumeSurfaceAreaStarter_urduKey.tex
%   tier 3 vocabulary   green   RGB 0 128 0
%   English text        black   RGB 0 0 0
%   Urdu text           purple  RGB 112 48 160
%   layout macros       version 3
%
% Please don't edit any information in this comment block - the server
% compares it against the current settings and rebuilds this file when
% they differ, so changes here would be overwritten. However, feel free
% to make updates to the LaTeX below if you want to edit it.
% ================================================================

\documentclass{beamer}
% ================================================================
% Copied from the English sheet this was translated from, so that
% anything it sets up for itself still works here
% ================================================================
\usepackage{tikz}
\usepackage{xstring}
\usepackage{amsmath}
\usetikzlibrary{calc}

% ================================================================
% Answer-overlay helpers
% ================================================================
\newcommand{\ablank}[1]{%
  \alt<2>{\textcolor{red}{#1}}{\underline{\phantom{#1}}}%
}
\newcommand{\ashow}[1]{\uncover<2->{\textcolor{red}{\small #1}}}
\newcommand{\ashowq}[1]{\alt<2->{\textcolor{red}{\small #1}}{?}}

% Answers drawn straight on to a TikZ diagram, revealed on overlay 2 like \ashow.
% Space is reserved on overlay 1, so the diagram does not move between slides.
\usetikzlibrary{overlay-beamer-styles}
\tikzset{
  ans/.style     = {red, thick, visible on=<2->},
  ansfill/.style = {red!15, visible on=<2->},
  anslab/.style  = {red, font=\tiny, visible on=<2->},
}

\title{Starter}
\date{}

\newcommand{\cuboid}[4]{
\begin{scope}[shift={(#1,#2)}]
  % Draw the cuboid
  \draw (0,0) rectangle (2,1);
  \draw (0.6,0.6) rectangle (2.6,1.6);
  \draw (0,0) -- (0.6,0.6);
  \draw (2,0) -- (2.6,0.6);
  \draw (2,1) -- (2.6,1.6);
  \draw (0,1) -- (0.6,1.6);

  % Labels (letter above, dimensions below)
  \node at (1,-0.4) {\tiny \bf{#3}};
  \node at (1,-1.2) {\tiny #4};

  % ----- Parse dimension string #4 (format: width×height×depth) -----
  \StrBefore{#4}{×}[\widthstr]                % first number
  \StrBehind{#4}{×}[\temp]                    % rest = height×depth
  \StrBefore{\temp}{×}[\heightstr]             % second number
  \StrBehind{\temp}{×}[\depthstr]               % third number

  % ----- Width dimension (bottom front edge) -----
  \draw (0,0) -- (0,-0.6);                     % left extension
  \draw (2,0) -- (2,-0.6);                      % right extension
  \draw (0,-0.8) -- (2,-0.8);                   % dimension line
  % ticks at ends (vertical)
  \draw (0,-0.75) -- (0,-0.85);
  \draw (2,-0.75) -- (2,-0.85);
  % label
  \node at (1,-0.8) [fill=white, inner sep=1pt] {\tiny \widthstr};

  % ----- Height dimension (left front edge) -----
  \draw (0,0) -- (-0.6,0);                      % lower extension
  \draw (0,1) -- (-0.6,1);                       % upper extension
  \draw (-0.8,0) -- (-0.8,1);                    % dimension line
  % ticks at ends (horizontal)
  \draw (-0.75,0) -- (-0.85,0);
  \draw (-0.75,1) -- (-0.85,1);
  % label
  \node at (-0.8,0.5) [fill=white, inner sep=1pt] {\tiny \heightstr};

    % ----- Depth dimension (bottom right edge) -----
    \coordinate (B) at (2,0);
    \coordinate (F) at (2.6,0.6);
    \coordinate (dir) at ($(F)-(B)$);              % direction of edge (0.6,0.6)
    % outward offset perpendicular to edge (along (1,-1) direction)
    \coordinate (offset) at (0.3,-0.3);
    % shorten factor: move endpoints inward along dir
    \def\shorten{0.1}
    \coordinate (start) at ($(B)+(offset)+\shorten*(dir)$);
    \coordinate (end)   at ($(F)+(offset)-\shorten*(dir)$);
    % dimension line
    \draw (start) -- (end);
    % ticks perpendicular to dimension line (along direction of edge)
    \coordinate (tickDir) at ($(0,0)!0.1!(dir)$);        % vector for tick length
\coordinate (start) at ($(B)+(offset)+(0,0)!\shorten!(dir)$);
\coordinate (end)   at ($(F)+(offset)-(0,0)!\shorten!(dir)$);;
    % label
    \node at ($(start)!0.5!(end)$) [fill=white, inner sep=1pt] {\tiny \depthstr};
\end{scope}
}

% Characters this language's font has not got are borrowed from another,
% rather than being silently dropped - see docs/translated-sheet-latex.md
\directlua{%
  if luaotfload and luaotfload.add_fallback then
    luaotfload.add_fallback("eallatinfallback",
      { "Noto Serif:mode=harf;" })
  else
    tex.error("this luaotfload is too old to register a fallback font")
  end
}
\usepackage[english,bidi=basic]{babel}
\babelprovide[import]{urdu}
\babelfont[urdu]{rm}[Renderer=Harfbuzz, BoldFont={Noto Nastaliq Urdu}, ItalicFont={Noto Nastaliq Urdu}, BoldItalicFont={Noto Nastaliq Urdu}, RawFeature={fallback=eallatinfallback}]{Noto Nastaliq Urdu}
\babelfont[urdu]{sf}[Renderer=Harfbuzz, BoldFont={Noto Nastaliq Urdu}, ItalicFont={Noto Nastaliq Urdu}, BoldItalicFont={Noto Nastaliq Urdu}, RawFeature={fallback=eallatinfallback}]{Noto Nastaliq Urdu}
\newcommand{\ealtext}[1]{\foreignlanguage{urdu}{#1}}
\newcommand{\ealtextblock}[1]{{\raggedleft\foreignlanguage{urdu}{#1}\par}}
% ================================================================
% EAL layout helpers
% ================================================================
\usepackage{xcolor}

\definecolor{ealkeycolour}{RGB}{0,128,0}
\definecolor{ealenglish}{RGB}{0,0,0}
\definecolor{ealtwocolour}{RGB}{112,48,160}

\newsavebox{\ealboxen}
\newsavebox{\ealboxtr}
\newlength{\ealglosswd}
\newlength{\ealglossraise}
\setlength{\ealglossraise}{1.15em}

% a tier 3 word, in English and in translation alike
\newcommand{\ealkey}[1]{\textcolor{ealkeycolour}{#1}}
\newcommand{\ealkeytr}[1]{\textcolor{ealkeycolour}{\ealtext{#1}}}

% One translated word above one English word. Built from kernel
% primitives, never a tabular, which would break wherever the sheet
% loads array, booktabs or tabularx. The box is as wide as the wider of
% the two, so a long translation pushes its neighbours aside instead of
% overlapping them, and it claims its full height so a table row or a
% TikZ node makes room for it.
\newcommand{\ealgloss}[2]{%
  \sbox{\ealboxen}{\ealkey{#1}}%
  \sbox{\ealboxtr}{\small\textcolor{ealtwocolour}{\ealtext{#2}}}%
  \setlength{\ealglosswd}{\wd\ealboxen}%
  \ifdim\wd\ealboxtr>\ealglosswd\setlength{\ealglosswd}{\wd\ealboxtr}\fi
  \leavevmode
  \makebox[\ealglosswd][c]{%
    \makebox[0pt][c]{\usebox{\ealboxen}}%
    \makebox[0pt][c]{\raisebox{\ealglossraise}{\usebox{\ealboxtr}}}%
  }%
}

% opens up line spacing around a block of glosses, so the raised
% translations sit clear of the line above
\newenvironment{ealglossed}{\par\linespread{2.1}\selectfont\ignorespaces}{\par}

% a whole translated sentence above its English counterpart. block level,
% so it does not belong inside a tabular cell
\newcommand{\ealpara}[2]{%
  \par\addvspace{0.5\baselineskip}%
  {\color{ealtwocolour}\ealtextblock{#1}}%
  \penalty200
  {\color{ealenglish}#2\par}%
  \addvspace{0.5\baselineskip}%
}

\begin{document}

\usetikzlibrary{calc}

\providecommand{\cuboid}[4]{
\begin{scope}[shift={(#1,#2)}]
  % Draw the cuboid
  \draw (0,0) rectangle (2,1);
  \draw (0.6,0.6) rectangle (2.6,1.6);
  \draw (0,0) -- (0.6,0.6);
  \draw (2,0) -- (2.6,0.6);
  \draw (2,1) -- (2.6,1.6);
  \draw (0,1) -- (0.6,1.6);

  % Labels (letter above, dimensions below)
  \node at (1,-0.4) {\tiny \bf{#3}};
  \node at (1,-1.2) {\tiny #4};

  % ----- Parse dimension string #4 (format: width×height×depth) -----
  \StrBefore{#4}{×}[\widthstr]                % first number
  \StrBehind{#4}{×}[\temp]                    % rest = height×depth
  \StrBefore{\temp}{×}[\heightstr]             % second number
  \StrBehind{\temp}{×}[\depthstr]               % third number

  % ----- Width dimension (bottom front edge) -----
  \draw (0,0) -- (0,-0.6);                     % left extension
  \draw (2,0) -- (2,-0.6);                      % right extension
  \draw (0,-0.8) -- (2,-0.8);                   % dimension line
  % ticks at ends (vertical)
  \draw (0,-0.75) -- (0,-0.85);
  \draw (2,-0.75) -- (2,-0.85);
  % label
  \node at (1,-0.8) [fill=white, inner sep=1pt] {\tiny \widthstr};

  % ----- Height dimension (left front edge) -----
  \draw (0,0) -- (-0.6,0);                      % lower extension
  \draw (0,1) -- (-0.6,1);                       % upper extension
  \draw (-0.8,0) -- (-0.8,1);                    % dimension line
  % ticks at ends (horizontal)
  \draw (-0.75,0) -- (-0.85,0);
  \draw (-0.75,1) -- (-0.85,1);
  % label
  \node at (-0.8,0.5) [fill=white, inner sep=1pt] {\tiny \heightstr};

    % ----- Depth dimension (bottom right edge) -----
    \coordinate (B) at (2,0);
    \coordinate (F) at (2.6,0.6);
    \coordinate (dir) at ($(F)-(B)$);              % direction of edge (0.6,0.6)
    % outward offset perpendicular to edge (along (1,-1) direction)
    \coordinate (offset) at (0.3,-0.3);
    % shorten factor: move endpoints inward along dir
    \def\shorten{0.1}
    \coordinate (start) at ($(B)+(offset)+\shorten*(dir)$);
    \coordinate (end)   at ($(F)+(offset)-\shorten*(dir)$);
    % dimension line
    \draw (start) -- (end);
    % ticks perpendicular to dimension line (along direction of edge)
    \coordinate (tickDir) at ($(0,0)!0.1!(dir)$);        % vector for tick length
\coordinate (start) at ($(B)+(offset)+(0,0)!\shorten!(dir)$);
\coordinate (end)   at ($(F)+(offset)-(0,0)!\shorten!(dir)$);;
    % label
    \node at ($(start)!0.5!(end)$) [fill=white, inner sep=1pt] {\tiny \depthstr};
\end{scope}
}

\begin{frame}

\begin{columns}[T]

% LEFT COLUMN
\column{0.5\textwidth}

\textbf{1.}

\begin{tikzpicture}[scale=0.55]
\def\R{3}
\def\r{1.5}

\fill[gray!30] (0,0) circle (\R);
\fill[white] (0,0) circle (\r);

\draw (0,0) circle (\R);
\draw (0,0) circle (\r);

\draw (0,0) -- (\R,0);
\draw (0,0) -- (0,\r);

% FIXED: Use separate coordinates for midpoint calculation
\path (0,0) -- (\R,0) coordinate[midway] (midR);
\path (0,0) -- (0,\r) coordinate[midway] (midr);
\node at (midR) [yshift=-0.3cm] {$6$ cm};
\node at (midr) [xshift=+0.1cm] {$3$ cm};
\end{tikzpicture}

\vspace{0.2cm}
\begin{ealglossed}
Find the shaded \ealgloss{area}{رقبہ}.
\end{ealglossed}

\begin{ealglossed}
\uncover<2->{\small \ealgloss{Area}{رقبہ} = \(27\pi\text{ cm}^2\)}
\end{ealglossed}

\vspace{0.4cm}

\textbf{3.}

\begin{tikzpicture}[scale=0.45]

\draw (0,0) rectangle (3,2);
\node at (1.5,-0.4) {\small A};
\node at (1.5,2.5) {6};
\node at (-0.4,1) {4};

\draw (5,0) rectangle (9.5,3);
\node at (7.25,-0.4) {\small B};
\node at (7.25,3.5) {9};
\node at (4.6,1.5) {6};

\end{tikzpicture}

\begin{ealglossed}
Find the \ealgloss{scale factor}{پیمانہ عامل} \ealgloss{mapping}{نگاشت} \newline A → B.
\end{ealglossed}

\begin{ealglossed}
\uncover<2->{\small \ealgloss{Scale factor}{پیمانہ عامل} = \(1.5\)}
\end{ealglossed}

% RIGHT COLUMN
\column{0.5\textwidth}

\textbf{2.}

\begin{tikzpicture}[scale=0.45]
\def\R{3}
\def\r{1.5}
\def\angle{60}

% Shaded annular sector
\fill[gray!30] (0,0) -- (0:\R) arc (0:\angle:\R) -- cycle;
\fill[white] (0,0) -- (0:\r) arc (0:\angle:\r) -- cycle;

% Draw only the arcs that bound the shaded region (no full circles)
\draw (0:\R) arc (0:\angle:\R);
\draw (0:\r) arc (0:\angle:\r);

% Draw the two radii
\draw (0,0) -- (0:\R);
\draw (0,0) -- (\angle:\R);

% Labels
\node at ($(0,0)!0.5!(\R,0)$) [yshift=-0.3cm] {$6$ cm};
\node at ($(0,0)!0.5!(0,\r)$) [xshift=-0.3cm] {$3$ cm};
\node at (\angle/2:2.2) {$60^\circ$};

% Add question text to the right of the diagram
\node[anchor=west] at (3.5,1.5) {Find the shaded \ealgloss{area}{رقبہ}.};
\end{tikzpicture}

\uncover<2->{\small \(\frac{1}{6}\times 27\pi = 4.5\pi\text{ cm}^2\)}

\vspace{0.2cm}

\textbf{4.}
\begin{center}
 \begin{tikzpicture}[scale=0.5]
    \cuboid{0}{0}{X}{4×6×8}
\end{tikzpicture}
\end{center}

\begin{tikzpicture}[scale=0.5]
  \cuboid{0}{-2}{A}{8×12×20}
  \cuboid{5}{-2}{B}{6×9×12}
  \cuboid{0}{-6}{C}{5×7.5×10}
  \cuboid{5}{-6}{D}{3×4.5×6}
\end{tikzpicture}

\begin{ealglossed}
Which are \ealgloss{similar}{مشابہ} to $X$?
\end{ealglossed}

\begin{ealglossed}
\uncover<2->{\small \ealgloss{Similar}{مشابہ} shapes: B, C, D}
\end{ealglossed}

\end{columns}

\end{frame}

%------------------------------------------------
% ANSWERS SLIDE
%------------------------------------------------

\begin{frame}

\begin{columns}[T]

% LEFT COLUMN
\column{0.5\textwidth}

\textbf{1.}

\begin{tikzpicture}[scale=0.55]
\def\R{3}
\def\r{1.5}

\fill[gray!30] (0,0) circle (\R);
\fill[white] (0,0) circle (\r);

\draw (0,0) circle (\R);
\draw (0,0) circle (\r);

\draw (0,0) -- (\R,0);
\draw (0,0) -- (0,\r);

% FIXED: Use separate coordinates for midpoint calculation
\path (0,0) -- (\R,0) coordinate[midway] (midR);
\path (0,0) -- (0,\r) coordinate[midway] (midr);
\node at (midR) [yshift=-0.3cm] {$6$ cm};
\node at (midr) [xshift=0.1cm] {$3$ cm};
\end{tikzpicture}

\(\pi6^2 - \pi3^2 = 27\pi\ \text{cm}^2\)

\vspace{0.4cm}

\textbf{3.}

\begin{tikzpicture}[scale=0.45]

\draw (0,0) rectangle (3,2);
\node at (1.5,-0.4) {\small A};
\node at (1.5,2.5) {6};
\node at (-0.4,1) {4};

\draw (5,0) rectangle (9.5,3);
\node at (7.25,-0.4) {\small B};
\node at (7.25,3.5) {9};
\node at (4.6,1.5) {6};

\end{tikzpicture}

\begin{ealglossed}
\ealgloss{Scale factor}{پیمانہ عامل} = \(1.5\)
\end{ealglossed}

% RIGHT COLUMN
\column{0.5\textwidth}

\textbf{2.}

\begin{tikzpicture}[scale=0.45]
\def\R{3}
\def\r{1.5}
\def\angle{60}

% Shaded annular sector
\fill[gray!30] (0,0) -- (0:\R) arc (0:\angle:\R) -- cycle;
\fill[white] (0,0) -- (0:\r) arc (0:\angle:\r) -- cycle;

% Draw only the arcs that bound the shaded region (no full circles)
\draw (0:\R) arc (0:\angle:\R);
\draw (0:\r) arc (0:\angle:\r);

% Draw the two radii
\draw (0,0) -- (0:\R);
\draw (0,0) -- (\angle:\R);

% Labels
\node at ($(0,0)!0.5!(\R,0)$) [yshift=-0.3cm] {$6$ cm};
\node at ($(0,0)!0.5!(0,\r)$) [xshift=-0.3cm] {$3$ cm};
\node at (\angle/2:2.2) {$60^\circ$};

% Add answer text to the right of the diagram
\node[anchor=west] at (3.5,1.5) {\(\frac{60}{360} \times 27\pi = \frac{1}{6}\times27\pi\)};
\node[anchor=west] at (3.5,0.8) {\(= 4.5\pi\ \text{cm}^2\)};
\end{tikzpicture}

\vspace{0.2cm}

\textbf{4.}

\begin{center}
\begin{tikzpicture}[scale=0.5]
    \cuboid{0}{0}{X}{4×6×8}
\end{tikzpicture}
\end{center}

\begin{tikzpicture}[scale=0.5]
  \cuboid{0}{-2}{A}{8×12×20}
  \cuboid{5}{-2}{B}{6×9×12}
  \cuboid{0}{-6}{C}{5×7.5×10}
  \cuboid{5}{-6}{D}{3×4.5×6}
\end{tikzpicture}

\begin{ealglossed}
\ealgloss{Similar}{مشابہ} shapes: B, C, D
\end{ealglossed}

\end{columns}

\end{frame}

\end{document}